% !TeX root = ../drafts/10-ethereum.tex
\section{Ethereum use cases}\label{sec:ethereum}

An aggregate proof certifies signature claims and the correct encoding of blob data. One guest program establishes these claims by checking raw signatures, checking a blob matrix, or recursively verifying child proofs of the same program.

\subsection{Recursive aggregation}\label{sec:recursion}\label{sec:xmss-agg}\label{sec:sphincs-agg}

A signature claim identifies what was signed and by whom. For XMSS~\cite{rfc8391}, it is a public key, an epoch, and a message; the keys claiming one $(\text{epoch},\text{message})$ pair form one group, so an epoch signed at under several messages contributes one group per message. For SPHINCS+~\cite{sphincs}, it is a $(\text{key},\text{message})$ pair.

A node may publish any subset of the signature claims supported by its raw signatures and verified children. The default is their sorted, deduplicated union. The essential condition is coverage: every published claim must be justified, either directly at this node or recursively through a child.

The DA claim is a sorted, duplicate-free list of at most $16$ LeanDA roots, possibly empty. Each root is paired with the hash of its membership test vector $\dual$ (\S\ref{sec:da-membership}). 

\subsubsection{The public input}\label{sec:rec-stmt}

The proof carries the signature claims and DA root list. Their digests bind these variable-length lists into a fixed-size statement. The VM's public input is BLAKE2s of the following fields, in order:
\begin{center}
\begin{tabular}{@{}rl@{}}
\hline
$32$ bytes & Fiat-Shamir IV\\
$32$ bytes & signer-set digest\\
$32$ bytes & DA list digest\\
$32$ bytes & credit digest\\
$16$ bytes & omitted SPHINCS+ count $\gen^n$\\
$(\kbc+5)\times24$ bytes & deferred bytecode point and value\\
$30\times24$ bytes & deferred matrix point and two values\\
\hline
\end{tabular}
\end{center}
The deferred claims let recursion postpone expensive polynomial evaluations until the final verification (\S\ref{sec:deferred-claims}).

The DA list digest is BLAKE2s of the concatenated pairs $(\daroot,\davechash)$, each a $32$-byte DA root $\daroot$ followed by a $32$-byte hash $\davechash$ of the membership test vector $\dual$.

The signer-set digest is BLAKE2s of the following, where each cell is $16$ bytes and a count is encoded as $\gen^{\text{count}}$ (8 bytes, zero-padded to one cell):
\begin{center}
\begin{tabular}{@{}ll@{}}
\hline
block $0$ & XMSS group count, SPHINCS+ claim count, SPHINCS+ list digest\\
per XMSS group & epoch, key count, $32$-byte message\\
& key list digest, two zero cells\\
\hline
\end{tabular}
\end{center}
Each inner list digest is BLAKE2s of the list alone: $32$ bytes per XMSS key or $64$ bytes per SPHINCS+ $(\text{key},\text{message})$ pair.

The credit digest names, for each published SPHINCS+ claim, every payee that did work for it, and how much. It is BLAKE2s of four $32$-byte digests: of the set sizes, of the entries' payees, of the entries' weights, and of the payee table. The payee table lists $20$-byte addresses, each zero-padded to two cells; entry $0$ is the payee of whoever proves this node, and the others are sorted and distinct from it. The honest prover lists only the payees its credits name; the verifier does not require that. Each published SPHINCS+ claim, in list order, has a non-empty set of entries $(\text{payee}, \text{weight})$, strictly increasing by payee. The size list holds one cell per claim, $\gen^m$ for a set of $m$; the payee list holds the entries' payees one after another, $\gen^i$ for payee $i$; the weight list, parallel to it, $\gen^W$ for weight $W$. Weights count work units: $L = 1024$ for a raw signature proven in a leaf, and $M = 77L$ for one merge of two proofs, about the ratio of their costs. Each list hash binds its own length, and the sizes split the other two into claims. The signer-set digest does not change, so it remains a function of the claims alone. The omitted count is the number $n$ of SPHINCS+ extra slots that hold claims this node covers and does not publish (\S\ref{sec:rec-coverage}), as $\gen^n$ in the low lane of one cell.

\subsubsection{Checking coverage}\label{sec:rec-coverage}

Every published claim needs support from a direct check or a recursively verified child. The guest builds a table with a region per XMSS (epoch, message) group and separate regions for SPHINCS+ and DA. Each region starts with the claims the parent publishes, followed by extra slots for duplicates and omitted claims.

Each occurrence of a supporting claim needs its own slot. If two children prove the same published claim, one occurrence uses its published slot and the other uses an extra slot. Every occurrence of a claim the parent omits uses an extra slot. These claims are still verified; only the published prefixes enter the parent's list digests.

The SPHINCS+ extra slots come in two parts, linked copies and then omitted claims. A linked copy is a second occurrence of a published claim: the prover hints the index of that claim's published slot, the guest range-checks it against the published count and copies the claim's four cells from that slot, so the copy is exactly the published claim. The omitted slots hold claims the prover supplies, and the statement publishes their count, per node: each child's count enters only the child statement the parent rebuilds.

Each omitted claim is shown absent from the published list, so no copy of a published claim can sit among the omitted slots, and every occurrence of a published SPHINCS+ claim a node holds is either in the claim's published slot or linked to it. The published list is sorted by claim, so the guest first puts it in key order. Each claim $c$ gets a $16$-bit key $k(c)$. With $\ell \in K$ the first $8$ bytes of $\text{BLAKE2s}(c)$, read as a lane, and $N(x) = x^{1+2^{16}+2^{32}+2^{48}}$ the norm from $K$ to its $2^{16}$-element subfield, $k(c)$ is the $e < 2^{16}-1$ with $N(\gen)^e = N(\ell)$, and $0$ when $\ell = 0$. $N$ maps $K^*$ evenly onto the subfield's $2^{16}-1$ nonzero elements, and $N(\gen)$ generates them, so the key is uniform and has one value. The prover hints $x = \gen^e$ and a root $u$. The guest range-checks $e < 2^{16}-1$, requires $u \in K$ and $u \ne 0$, and checks $x\,u = \ell\,u^{2^{16}}$: then $x/\ell = u^{2^{16}-1}$ lies in the kernel of $N$, so $N(x) = N(\ell)$. Requiring $u \in K$ matters: in $F_{192}$ the $(2^{16}-1)$-th powers meet $K^*$ in a subgroup three times the kernel, which would leave three keys. For $\ell = 0$ it requires $x = 1$ instead. The prover hints, for each position of the key order, which published claim sits there. The guest writes a distinct mark at that claim's index, which makes the hint a permutation, and requires the keys not to decrease along it. Sentinels $\gen^0$ and $\gen^{2^{16}}$ bound the order at both ends, the keys being stored as $\gen^{k+1}$. For an omitted claim $c'$, the prover hints $p$, the number of published keys below $k(c')$, and $r$, the number equal to it. The guest checks that the key at position $p-1$ is smaller and the one at $p+r$ larger, both strictly, that the $r$ between equal $k(c')$, and that each of them differs from $c'$ in a hinted cell. Since keys never decrease along the order, every published claim with key $k(c')$ is among those $r$, so $c'$ is none of them. Collisions cost only the $r$ comparisons, never soundness; $r$ is $0$ except with probability about $n/2^{16}$. A node that omits nothing skips the pass. Otherwise it costs one compression and about $80$ VM steps per published claim, and about as much per omitted claim.

The prover supplies the extra claims and assigns each supporting occurrence to a slot. The guest checks that they match, confines the assignment to the correct region, and writes a distinct counter value there. Write-once memory prevents slot reuse. Requiring as many writes as slots ensures that every slot is covered, including every published claim.

A SPHINCS+ slot also records its occurrence's credit entries, next to its coverage mark and therefore exactly once: where their payees and weights are, how many, and the map from their payee indices into this node's table. A raw signature records $\{(\gen^0, \gen^L)\}$, this node with one leaf's work, with the identity map. A child's claim records the child's entries for that claim, in the child's indices, with a map from the child's payee table to this node's: one hinted entry per child payee, either $\gen^k$, which the guest checks by comparing the two addresses, or $0$ for a payee this node dropped. The guest reads the child's table and its three credit lists from hints and hashes them into the child's credit digest, which enters the child statement it rebuilds, so none can be altered. It range-checks every payee of the child's sets against the child's table, since it indexes the map with them. The child's weights need no check: the child proved each one.

The node's own sets are hinted too, and the guest proves each one from the claim's occurrences: its published slot and every copy linked to it, which, with the omitted claims shown absent, are all the occurrences it holds. Every element of every occurrence, mapped, is a contribution to the entry of its claim's set with that payee, at a hinted position. Each entry's weight is the \emph{maximum} of its contributions, so work that reaches a node by two paths counts once. Each entry names, by a hinted id, the contribution that attains its maximum, and that contribution must equal it; every other contribution must not exceed it, by a hinted difference in two $16$-bit limbs, exact because all weights and differences stay far below the order of $\gen$. Every contribution records the entry it joined, and a final pass checks that each entry's designated contribution joined that entry, which also rules out an entry nobody contributed to. The sets' payees are strictly increasing, so one set has one encoding.

A fold of $m \ge 2$ children adds its own share of the merges, $f = \lfloor (m-1) M / n \rfloor$ for its $n$ published SPHINCS+ claims, to its own entry of every claim: a zero-weight contribution makes that entry exist, and its published weight is its maximum plus $f$. The prover hints $f$ in two $16$-bit limbs and the remainder $r < n$, and the guest checks $\gen^{fn}\gen^r = \gen^{(m-1)M}$, taking $\gen^{fn}$ by square-and-multiply over the bits of $n$; since $fn + r < 2^{48}$, the exponents are equal as integers. A one-child fold earns nothing. The children of a fold must be distinct proofs: every ordered pair differs in a hinted cell of their statements. A dropped payee maps to $0$, which is no power of $\gen$: a set holding it is rejected by whoever reads the list, the parent or the final verifier, so a node keeps only the payees its published claims name. The payee is part of the public input, which is absorbed before the first commitment, and the guest computes the credit digest itself, so a proof or a prover state made under one payee cannot be reused under another.

A claim's credits are therefore every node that verified it raw or merged it, along the paths of folds on which the claim stayed published, each with its heaviest path's work. A node that omits a claim drops the credits below it: if the claim returns later through another proof, only that proof's work comes with it. No node can drop a worker of a claim it publishes, add one, move one to another claim, or change a weight. Payout is left to whoever reads the sets: a transaction's budget can be split by weight. Two known limits of the maximum: a payee that did different work on two paths is paid its heaviest path only, and a merge of two distinct but overlapping aggregates still earns its share.

\subsubsection{Deferred evaluation claims}\label{sec:deferred-claims}

Verifying a child requires evaluations of three fixed polynomials: the bytecode multilinear, in $\kbc+4$ variables (\S\ref{sec:e2e-bc}), and flock's R1CS matrix multilinears $\widetilde A_0,\widetilde B_0$, each in $28$ variables (Annex~\ref{annex:flock}). Computing them inside the guest would be very expensive. Instead, the guest carries evaluation claims in its statement. These expensive computations are moved 'outside of the snark'.

A node combines the claims carried by its children with the new claims produced while verifying them. Two sumchecks reduce these to three evaluations: one of the bytecode polynomial and one of each matrix polynomial, batched. The Fiat-Shamir for batching and sumchecks is seeded by all child and running claims.

\subsection{Data availability: LeanDA}\label{sec:leanda}

Erasure coding makes data recoverable from a subset of its encoded symbols. Sampling checks whether those symbols can be retrieved, but does not establish that they form a valid encoding. An approach using leanVM was introduced in \href{https://frisitano.github.io/slides/presentations/leanvm-fold-in-circuit-das-2026-05-08/index.html}{Risitano's presentation} and developed in LeanDA~\cite{LeanDA} by Meng, Wagner, Kadianakis and Risitano: commit to an encoded matrix and prove every row is a codeword, then sample authenticated cells. A formal security proof can be found in \cite{LeanDASecurity}.

Our construction uses BLAKE2s and the additive Reed-Solomon encoder over $\K$ already used by the PCS, with a different codeword membership test from LeanDA's (\S\ref{sec:da-membership}). The polynomial basis, transform, and membership lemmas are collected in Annex~\ref{app:ntt}.


\subsubsection{The commitment}\label{sec:da-commit}

A blob is a polynomial in $\K[X]$ of degree less than $\kdim$, represented by its $\kdim$ evaluations on $\subsp_{\log_2\kdim}$. Its Reed-Solomon encoding is the vector of evaluations on the larger domain $\dom=\subsp_{\log_2\blen}$, where $\blen=2\kdim$. The first $\kdim$ symbols are the original blob, so the encoding is systematic; the remaining symbols provide redundancy. Any $\kdim$ encoded evaluations recover the blob.

One commitment covers a nonempty matrix of $\nblob$ encoded blobs, one per row.

Split each encoded row into cells of $\cellsz$ symbols. The guest fixes these dimensions, matching the byte sizes of an \href{https://eips.ethereum.org/EIPS/eip-4844}{EIP-4844} blob and an \href{https://eips.ethereum.org/EIPS/eip-7594}{EIP-7594} sampling cell:
\begin{center}
\begin{tabular}{@{}rll@{}}
\hline
$\kdim$ & $2^{14}$ symbols & $128$ KiB of payload per row\\
$\blen$ & $2^{15}$ symbols & $256$ KiB per encoded row\\
$\cellsz$ & $2^8$ symbols & $2$ KiB per cell\\
$\blen/\cellsz$ & $128$ cells & any $64$ reconstruct a row\\
$\nblob$ & $1$ to $1024$ & dynamic parameter\\
\hline
\end{tabular}
\end{center}
Zero codewords pad the row count to the next power of two. The root binds this padded matrix.

Let $H$ denote BLAKE2s and $e_{i,j}$ the hash of cell $j$ in row $i$. The same cell digests feed two branches:
\begin{itemize}[itemsep=1pt,topsep=3pt]
\item \textbf{Row branch.} Hash the first $\kdim/\cellsz$ cell digests of each row into $r_i$, then Merkle-commit these row digests into $\darow$.
\item \textbf{Column branch.} Merkle-commit each column of cell digests into $C_j$, then Merkle-commit the column roots into $\dacol$.
\end{itemize}
The matrix root is $\daroot=H(\darow,\dacol)$ (Figure~\ref{fig:da-hashing}). See \cite{LeanDA} for details.

\begin{figure}[htbp]
\centering
\begin{tikzpicture}[
  font=\footnotesize,
  >={Stealth[length=1.5mm]},
  edge/.style={semithick,draw=black!65},
  rowedge/.style={semithick,draw=blue!65!black},
  digest/.style={circle,fill=black!65,inner sep=1.5pt},
  rowdigest/.style={circle,fill=blue!65!black,inner sep=1.5pt}
]
\node[anchor=west,font=\small\bfseries] at (0,1.3) {Encoded blobs: one row per blob, $128$ cells per row};
\node[blue!65!black] at (2,0.75) {First $64$ cells};
\node at (6,0.75) {Remaining $64$ cells};
\foreach \j/\label in {0/0,1/1,2/\cdots,3/63,4/64,5/\cdots,6/126,7/127} {
  \node at (\j+0.5,0.25) {$\label$};
}
\foreach \i in {0,1,2,3} {
  \fill[blue!8] (0,-\i*0.65) rectangle (4,-\i*0.65-0.65);
  \foreach \j/\label in {0/0,1/1,2/\cdots,3/63,4/64,5/\cdots,6/126,7/127} {
    \draw[black!35] (\j,-\i*0.65) rectangle (\j+1,-\i*0.65-0.65);
    \ifnum\j=2
      \node at (\j+0.5,-\i*0.65-0.325) {$\cdots$};
    \else\ifnum\j=5
      \node at (\j+0.5,-\i*0.65-0.325) {$\cdots$};
    \else
      \node at (\j+0.5,-\i*0.65-0.325) {$e_{\i,\label}$};
    \fi\fi
  }
  \node[rowdigest,label=above:{$r_\i$}] (row\i) at (-0.9,-\i*0.65-0.325) {};
  \draw[rowedge,->] (0,-\i*0.65-0.325) -- (row\i);
}
\draw[blue!65!black,thick] (0,0) rectangle (4,-2.6);
\draw[orange!85!black,thick] (1,0) rectangle (2,-2.6);
\node[align=center,blue!65!black] at (-1.8,0.75) {$r_i=H(e_{i,0},\ldots,e_{i,63})$};
\node[rowdigest] (row01) at (-1.8,-0.65) {};
\node[rowdigest] (row23) at (-1.8,-1.95) {};
\node[rowdigest,label=above:{$\darow$}] (rowroot) at (-2.9,-1.3) {};
\draw[rowedge] (row0) -- (row01) -- (row1);
\draw[rowedge] (row2) -- (row23) -- (row3);
\draw[rowedge] (row01) -- (rowroot) -- (row23);
\node[align=center] at (5,-3.05) {Each rectangle: one $2$ KiB cell, labelled by\\its digest $e_{i,j}=H(\text{cell}_{i,j})$.};

\draw[orange!85!black,dashed,->] (1.5,-2.6) -- (1.5,-3.65);
\node[anchor=west,font=\small\bfseries] at (-0.3,-4) {Expand column $1$};
\node[anchor=west] at (3.1,-4) {The same tree is built for every column.};
\foreach \i in {0,1,2,3} {
  \node[draw=orange!85!black,fill=orange!8,minimum width=1.05cm,minimum height=0.45cm] (colleaf\i) at (\i*1.6,-4.65) {$e_{\i,1}$};
}
\node[digest] (col01) at (0.8,-5.4) {};
\node[digest] (col23) at (4,-5.4) {};
\node[digest,label=below right:{$C_1$}] (col1) at (2.4,-6.15) {};
\draw[edge] (colleaf0) -- (col01) -- (colleaf1);
\draw[edge] (colleaf2) -- (col23) -- (colleaf3);
\draw[edge] (col01) -- (col1) -- (col23);

\node[digest,label=above:{$C_0$}] (col0) at (0.4,-6.15) {};
\node[digest,label=above:{$C_{126}$}] (col126) at (6.4,-6.15) {};
\node[digest,label=above:{$C_{127}$}] (col127) at (8.4,-6.15) {};
\node at (4.4,-6.15) {$\cdots$};
\node[digest] (cols01) at (1.4,-6.9) {};
\node[digest] (cols126127) at (7.4,-6.9) {};
\draw[edge] (col0) -- (cols01) -- (col1);
\draw[edge] (col126) -- (cols126127) -- (col127);
\node[digest,label=below:{$\dacol$}] (colroot) at (4.4,-8.2) {};
\draw[edge,dashed] (cols01) -- (colroot) -- (cols126127);
\node[fill=white,inner sep=3pt] at (4.4,-7.35) {Merkle tree over all $128$ column roots};

\node[draw,rounded corners=2pt,fill=black!4,inner sep=5pt] (root) at (0.5,-9.35) {$\daroot=H(\darow,\dacol)$};
\draw[rowedge,->] (rowroot) -- (-2.9,-9.35) -- (root.west);
\draw[edge,->] (colroot) -- (4.9,-8.2) |- (root.east);
\end{tikzpicture}
\caption{Example of encoding of 4 blobs.}
\label{fig:da-hashing}
\end{figure}

\subsubsection{Codeword membership}\label{sec:da-membership}\label{sec:da-circuit}

A row is valid exactly when it consists of evaluations of a polynomial of degree less than $\kdim$. On our additive domain at rate $1/2$, such rows are orthogonal to every codeword, and this condition also characterizes them: the code is self-dual (Corollary~\ref{cor:rs-dual}). This gives an efficient membership test.

The Fiat-Shamir challenges $\log_2\kdim$ in $\E$ are derived from the matrix root $\daroot$. These define the polynomial $\dual$ of \eqref{eq:dual-product} and its evaluation vector $\davec=(\dual(x))_{x\in\dom}$. Let $\davechash=H(\davec)$, serializing each entry as its three little-endian $64$-bit limbs, in domain order.

The guest recomputes $\daroot$ from the encoded rows, receives $\davec$ as a hinted vector, checks its hash against $\davechash$, and uses it to check membership:
\[
  \inner{\davec}{w_i}
  =\sum_{x\in\dom}\dual(x)w_i(x)
  \;\qeq\;0.
\]
for every row $w_i$. The vector is shared across all rows and hashed once. It contains $32768$ entries, or $768$ KiB.

Recursive nodes preserve the pair $(\daroot,\davechash)$ for each retained claim. The final native verifier recomputes $\davechash$ from every retained root (outside of the snark).

Every valid row passes. Under uniform challenges, a fixed invalid matrix passes with probability at most $\log_2\kdim/|\E|=14/2^{192}$ (Proposition~\ref{prop:rs-membership}).
